% -*- coding: utf-8 -*-
% =============================================================================
% header.tex —— 自学案：导言区 + 宏定义 + 标题区 + 双栏开始
% 架构: student.tex/teacher.tex -> main.tex -> [header.tex, sec*.tex, footer.tex]
% 参数: 全部版面与文案参数来自 config-class.tex / config.tex
%       （由 scripts/gen_config.py 依据 config.yaml 自动生成，勿手改）。
% =============================================================================

% ---- 文档类选项（字号等，须在 \documentclass 之前注入） ----
% Auto-generated by scripts/gen_config.py from config.yaml
% 请勿手改本文件；修改 config.yaml 后重新运行 gen_config.py。
\PassOptionsToClass{12pt}{ctexart}

\documentclass{ctexart}

% ---- 版面几何与文本参数（config.yaml 自动生成） ----
% Auto-generated by scripts/gen_config.py from config.yaml
% 请勿手改本文件；修改 config.yaml 后重新运行 gen_config.py。

% --- 纸张与版面几何 ---
\usepackage[
  paperwidth=285mm,
  paperheight=420mm,
  landscape,
  top=1in,
  bottom=1in,
  left=1in,
  right=1in,
  headheight=25pt
]{geometry}

% --- 文本参数 (\pd<Key>) ---
\providecommand{\pdColumns}{2}
\providecommand{\pdMainTitle}{章节主标题（例如：群论）}
\providecommand{\pdTeacherHeadLeft}{【教师用书·参考解答与教学要点】}


% ---- 教师版版面覆盖：A4 竖版、单栏 ----
% 学生版沿用 config.yaml 的大8开横版双栏；教师版供案头翻阅，固定为 A4 竖版单栏。
% \geometry 在导言区可重复调用并覆盖先前设置；\pdColumns=1 时 \pdBeginColumns 不产生 multicols 环境。
\ifdefined\TeacherVersion
  \geometry{a4paper, portrait, top=2.5cm, bottom=2.5cm, left=2.5cm, right=2.5cm, headheight=25pt}
  \renewcommand{\pdColumns}{1}
\fi

% ---- 栏环境起止（统一由此处负责，main.tex/footer.tex 只调用宏） ----
% multicols 不会因参数为 1 自动退化为单栏（会警告并强制改为双栏），
% 因此栏数为 1 时不得进入 multicols 环境。
\newcommand{\pdBeginColumns}{%
  \ifnum\pdColumns>1 \begin{multicols}{\pdColumns}\fi
}
\newcommand{\pdEndColumns}{%
  \ifnum\pdColumns>1 \end{multicols}\fi
}

% -----------------------------------------------------------------------------
% 核心宏包加载
% -----------------------------------------------------------------------------
\usepackage{amsmath,amssymb,amsthm,bm,mathtools}
\usepackage{fancyhdr}
\usepackage{multicol}
\usepackage{enumitem}
\usepackage{graphicx}
\usepackage{xcolor}
\usepackage{comment}
\usepackage[framemethod=TikZ]{mdframed}
\usepackage{tikz}
\usetikzlibrary{arrows.meta, positioning, calc, shapes.geometric}
\usepackage{exam-zh-choices}  % 选择题选项自动排版宏包
\usepackage{hyperref}

\hypersetup{
  colorlinks=true,
  linkcolor=blue!75!black,
  citecolor=blue!75!black,
  urlcolor=blue!75!black
}

% -----------------------------------------------------------------------------
% 统计计数器（编译后输出环境使用摘要，便于验收与排错）
% -----------------------------------------------------------------------------
\newcounter{pdSolutionCount}
\newcounter{pdNoteCount}
\newcounter{pdFillinCount}
\newlength{\fillinwidth}

% -----------------------------------------------------------------------------
% 双版本环境控制宏 (Single-source Dual-output)
% 说明: solution/teacherNote/fillin/ansspace 在两版中的行为在此处统一定义，
%       切片文件 (sec*.tex) 只需使用同一套命令，无需关心版本差异。
% -----------------------------------------------------------------------------
\ifdefined\TeacherVersion
  % ============================ 教师版环境配置 ============================
  \pagestyle{fancy}
  \fancyhf{}
  \fancyhead[L]{\bfseries\color{blue!80!black}\pdTeacherHeadLeft}
  \fancyfoot[C]{\thepage\ of \pageref{LastPage}}
  \renewcommand{\headrulewidth}{0.4pt}
  \renewcommand{\footrulewidth}{0.4pt}

  % 题型标签：教师版显示【xx题】字样，便于把握设计意图
  \newcommand{\qtype}[1]{【#1】}

  % 题目状态标签 \practicemode{guided|unprompted}（可选，缺省 guided）：
  % unprompted 为无提示纯习题，教师版显示【无提示练习】标签，学生版隐藏。
  % 用 \providecommand 兼容，取值非法时发出编译警告。
  \providecommand{\practicemode}[1]{%
    \def\pdPmode{#1}%
    \def\pdPmodeU{unprompted}%
    \def\pdPmodeG{guided}%
    \ifx\pdPmode\pdPmodeU
      {\bfseries\color{teal!80!black}【无提示练习】}%
    \else
      \ifx\pdPmode\pdPmodeG\else
        \PackageWarning{problem-driven}{\protect\practicemode\space 取值只能是 guided 或 unprompted（收到“#1”）}%
      \fi
    \fi}

  % 参考解答环境（显示完整答案推导，并计数）
  \newenvironment{solution}[1][解]{%
    \stepcounter{pdSolutionCount}%
    \par\begingroup\color{blue!85!black}%
    \begin{proof}[\indent\bfseries\color{blue!85!black}【#1】]%
  }{%
    \end{proof}\endgroup\par
  }

  % 教师专用教学要点、易错警示与评分提示环境
  \newenvironment{teacherNote}[1][教学要点与易错分析]{%
    \stepcounter{pdNoteCount}%
    \par\vspace{0.6ex}%
    \begingroup\color{teal!80!black}\small
    \noindent\textbf{【#1】}\par\vspace{0.2ex}%
  }{%
    \endgroup\par\vspace{0.6ex}%
  }

  % 答题留白：教师版压缩为微小行距，消除多余空白
  \newcommand{\ansspace}[1]{\par\vspace{0.4ex}}

  % 填空宏：教师版在下划线居中呈现醒目的加粗蓝色答案；
  % 缺宽度或缺答案时发出编译警告
  \newcommand{\fillin}[2][]{%
    \stepcounter{pdFillinCount}%
    \if\relax\detokenize{#2}\relax
      \PackageWarning{problem-driven}{\protect\fillin\space 缺少预留宽度参数（第 \thepdFillinCount\space 处）}%
    \fi
    \if\relax\detokenize{#1}\relax
      \PackageWarning{problem-driven}{\protect\fillin\space 未提供参考答案（第 \thepdFillinCount\space 处）}%
      \underline{\hspace{#2}}%
    \else
      \ifmmode
        \settowidth{\fillinwidth}{$\bm{#1}$}%
        \ifdim\fillinwidth>#2
          \underline{\hspace{0.5em}{\color{blue!85!black}\bm{#1}}\hspace{0.5em}}%
        \else
          \underline{\makebox[#2][c]{$\color{blue!85!black}\bm{#1}$}}%
        \fi
      \else
        \settowidth{\fillinwidth}{\bfseries#1}%
        \ifdim\fillinwidth>#2
          \underline{\hspace{0.5em}\textbf{\color{blue!85!black}#1}\hspace{0.5em}}%
        \else
          \underline{\makebox[#2][c]{\bfseries\color{blue!85!black}#1}}%
        \fi
      \fi
    \fi
  }
\else
  % ============================ 学生版环境配置 ============================
  \pagestyle{fancy}
  \fancyhf{}
  \fancyhead[L]{姓名：\underline{\hspace{2.2cm}}\quad 班级：\underline{\hspace{2cm}}}
  \fancyfoot[C]{\thepage\ of \pageref{LastPage}}
  \renewcommand{\headrulewidth}{0.4pt}
  \renewcommand{\footrulewidth}{0.4pt}

  % 题型标签：学生版隐藏【xx题】字样（设计流程保留题型，但不明示）
  \newcommand{\qtype}[1]{}

  % 题目状态标签：学生版隐藏（unprompted 题的引导隐藏由切片写作纪律保证，
  % 并由 scripts/check_numbering.py 检查题块内不得含提示标记）
  \providecommand{\practicemode}[1]{}%

  % 彻底剥离答案与教师批注，同时计数以便编译日志中核对剥离数量
  % 注意: \specialcomment 默认会“包含”环境体；在组内将 \ProcessCutFile 重定义为
  % 空即恢复剥离效果（comment 宏包官方推荐用法），begin 行尾内容会被自动丢弃。
  \specialcomment{solution}{\stepcounter{pdSolutionCount}\begingroup\def\ProcessCutFile{}}{\endgroup}
  \specialcomment{teacherNote}{\stepcounter{pdNoteCount}\begingroup\def\ProcessCutFile{}}{\endgroup}

  % 答题留白：学生版产生实际书写所需的物理垂直空间
  \newcommand{\ansspace}[1]{\vspace{#1}}

  % 填空宏：学生版仅呈现预留长度的下划线空白
  \newcommand{\fillin}[2][]{%
    \stepcounter{pdFillinCount}%
    \if\relax\detokenize{#2}\relax
      \PackageWarning{problem-driven}{\protect\fillin\space 缺少预留宽度参数（第 \thepdFillinCount\space 处）}%
    \fi
    \underline{\hspace{#2}}%
  }
\fi

% 编译结束时的质量摘要（写入 .log 与终端，供 build.py 解析核对）
\AtEndDocument{%
  \typeout{[problem-driven] SUMMARY solution=\arabic{pdSolutionCount},
    teacherNote=\arabic{pdNoteCount}, fillin=\arabic{pdFillinCount}}%
  \ifdefined\TeacherVersion
    \ifnum\value{pdSolutionCount}=0
      \PackageWarning{problem-driven}{教师版未发现任何 solution 环境，请检查切片文件}%
    \fi
  \else
    \typeout{[problem-driven] 学生版已剥离 solution/teacherNote，数量为上述 SUMMARY}%
  \fi
}

% -----------------------------------------------------------------------------
% 微型知识点模块与列表规范
% -----------------------------------------------------------------------------
\newsavebox{\microbox}
\newenvironment{microknowledge}[1][核心定义与记号]{%
  \par\vspace{0.8ex}%
  \def\microtitle{#1}%
  \begin{lrbox}{\microbox}%
    \begin{minipage}{\dimexpr\linewidth-2\fboxsep-2\fboxrule\relax}%
      \textbf{【\microtitle】}\par\vspace{0.3ex}\small
}{%
    \end{minipage}%
  \end{lrbox}%
  \noindent\fbox{\usebox{\microbox}}\par\vspace{0.8ex}%
}

% -----------------------------------------------------------------------------
% 章节末结论框（虚线框）：只放本节由问题链自主探究推出的结论
% -----------------------------------------------------------------------------
\mdfdefinestyle{pdKnowledgeBoxStyle}{%
  linecolor=gray,
  linewidth=0.8pt,
  backgroundcolor=gray!4,
  tikzsetting={draw=gray, dash pattern=on 3pt off 2pt},
  roundcorner=0pt,
  innertopmargin=6pt,
  innerbottommargin=6pt,
  innerleftmargin=8pt,
  innerrightmargin=8pt,
  nobreak=true
}
\newenvironment{knowledgebox}[1][本节结论]{%
  \par\vspace{0.8ex}%
  \begin{mdframed}[style=pdKnowledgeBoxStyle]%
  \textbf{#1}\par\vspace{0.3ex}\small
}{%
  \end{mdframed}\par\vspace{0.8ex}%
}


% 列表样式统一设定 (一级题号连续递增；二级子题用 (a), (b)...)
\setlist[enumerate,1]{label=\textbf{\arabic*.}, leftmargin=*, itemsep=0.8em, topsep=0.4ex}
\setlist[enumerate,2]{label=(\alph*), leftmargin=*, itemsep=0.4em}

% -----------------------------------------------------------------------------
% 题链自动编号环境 probchain（推荐）
% 自建计数器跨环境续号，彻底摆脱手工维护 series=/resume= 的对位负担：
% 所有小节统一写作 \begin{probchain}...\end{probchain}，先后顺序天然决定题号。
% 旧写法 \begin{enumerate}[series=probchain] / [resume=probchain] 仍然兼容。
% -----------------------------------------------------------------------------
\newcounter{pdprob}
\newenvironment{probchain}{%
  \begin{enumerate}[start=\the\numexpr\value{pdprob}+1\relax]%
}{%
    \setcounter{pdprob}{\value{enumi}}%
  \end{enumerate}%
}
